% ============================================================================
% ML Defender (aRGus NDR) — v26-pre — DAY 255
% Contenido nuevo montado en orden de documento + correcciones in-place.
%
% "pre" = documenta el DIAGNÓSTICO de las cabezas y el PLAN de reparación;
%          los arreglos NO están ejecutados aún. Una v26 futura reportará
%          resultados de reparación medidos.
%
% INTEGRACIÓN:
%   A) Aplicar las 3 correcciones in-place del CHECKLIST (son ediciones, no añadidos).
%   B) Insertar los 4 bloques nuevos en los puntos indicados (o vía \input).
% Todo anclado a lo MEDIDO (DAY 255, rama diag/ml-heads).
% ============================================================================


% ╔══════════════════════════════════════════════════════════════════════════╗
% ║ CHECKLIST — 3 CORRECCIONES IN-PLACE (reemplazos, no añadidos)              ║
% ╠══════════════════════════════════════════════════════════════════════════╣
% ║ C1. §Implementation → "Sentinel value selection rationale": el v25 afirma  ║
% ║     que el centinela es "non-informative, not influencing classification". ║
% ║     Reemplazar ese párrafo por el bloque [C1] de abajo.                     ║
% ║                                                                            ║
% ║ C2. §Formal System Model → "Feature Extraction": reemplazar la oración     ║
% ║     "The sentinel components remain present in x_f by design: a missing    ║
% ║      feature must be in the vector to route deterministically to the left  ║
% ║      child of every split that evaluates it."  por el bloque [C2].         ║
% ║                                                                            ║
% ║ C3. §Limitations 10.2 "Incomplete Feature Set": reemplazar el cuerpo por   ║
% ║     el bloque [C3].                                                         ║
% ╚══════════════════════════════════════════════════════════════════════════╝


% ── [C1] — reemplaza el párrafo "Sentinel value selection rationale" ────────
\paragraph{Sentinel value selection rationale (revised, DAY~255).}
Split-domain analysis confirmed all Random Forest thresholds lie within
$[0.0, 5.1]$; $-9999.0\text{f}$ lies entirely outside this domain and routes
deterministically to the left child of any split that evaluates it. The
original design inferred from this that the sentinel is \emph{non-informative}.
A later code-level audit (DAY~255, \S\ref{sec:eval:deadfeatures}) refutes that
inference for any feature on which the trained forest actually splits:
deterministic left-routing is harmless only if the feature was \emph{also}
sentinel-valued during training. A decision-tree split is selected only where a
feature varies in the training set --- a constant feature yields zero
information gain and is never chosen --- so the presence of split nodes on a
feature is itself evidence that the feature varied at training time and is
constant only at inference. In that case the sentinel does not vanish from the
decision: it forces a fixed branch, a systematic train/serve skew.
$\texttt{quiet\_NaN()}$ was still correctly rejected (undefined cross-compiler
comparison); $0.5\text{f}$ was correctly rejected (within domain). The corrected
conclusion is narrower: the sentinel is a sound encoding of \emph{absence} only
for features absent in both training and serving, a condition
\S\ref{sec:eval:deadfeatures} shows the current models do not meet.

% ── [C2] — reemplaza la oración del §Formal System Model ────────────────────
The sentinel components remain present in $x_f$ by design, routing
deterministically to the left child of any split that evaluates them. We
initially treated this as classification-neutral; \S\ref{sec:eval:deadfeatures}
demonstrates by direct audit of the compiled forests that it is not neutral
whenever the forest was trained on a non-constant distribution of that feature,
and quantifies the resulting skew across all four heads.

% ── [C3] — reemplaza el cuerpo de §10.2 ────────────────────────────────────
\textbf{10.2 Incomplete Feature Set (measured impact, DAY~255).}
Of the 40 features in the ML~Defender contract, a subset carry the
$-9999.0\text{f}$ sentinel at inference because their computation is not
available in a network sniffer. \S\ref{sec:eval:deadfeatures} audits the
compiled forests and finds that all four heads allocate decision splits to
sentinel-served features --- 890 of 3{,}065 split nodes in total (29\%), rising
to 41.9\% for the Ransomware head, whose single most important feature
(\texttt{io\_intensity}, 24\% importance) measures disk I/O and is physically
inaccessible from network traffic. This is a train/serve skew, concurrent with
and independent of the distribution-transfer failure of
\S\ref{sec:eval:xgboost:ood}. It is a measured contributor to the low ML score
on the labeled Neris botnet (\S\ref{sec:eval:bias}); a remediation plan with
per-head cost is given in \S\ref{sec:eval:remediation}.


% ╔══════════════════════════════════════════════════════════════════════════╗
% ║ BLOQUE 1 — INSERTO DE ABSTRACT (opcional)                                  ║
% ║ Colocar tras el bloque "── NEW DAY 252 ──", antes de "ML Defender is       ║
% ║ released under the MIT license."                                           ║
% ╚══════════════════════════════════════════════════════════════════════════╝

% ── NEW DAY 255 ──────────────────────────────────────────────────────────
Finally (DAY~255), a code-level audit of the four embedded classifiers
establishes \emph{why} the ML score is blind on the labeled botnet independently
of distribution transfer: every head allocates decision splits to features
served as a constant absence-sentinel at inference (29\% of all split nodes;
41.9\% for the Ransomware head, whose highest-importance feature measures disk
I/O and cannot exist in a network capture). Because a forest never splits on a
constant feature, the existence of these splits proves the models were trained
on a feature space the deployment does not reproduce --- a train/serve skew
provable from the compiled trees alone. We report this as a measured,
reproducible failure mode and derive a per-head remediation plan whose principled
resolution for the Ransomware head is host-domain fusion with Wazuh rather than
further network feature engineering. The repairs are specified and costed here;
their measured outcomes are deferred to a subsequent revision.
% ──────────────────────────────────────────────────────────────────────────


% ╔══════════════════════════════════════════════════════════════════════════╗
% ║ BLOQUE 2 — NUEVA SUBSECCIÓN DE EVALUACIÓN                                  ║
% ║ Insertar al final de §Evaluation, tras §sec:eval:bias                      ║
% ╚══════════════════════════════════════════════════════════════════════════╝

\subsection{Dead-Feature Audit: A Train/Serve Skew Provable from the Compiled Trees}
\label{sec:eval:deadfeatures}

The per-lens analysis (\S\ref{sec:eval:bias}) established that on the labeled
Neris botnet the embedded ML score averages $0.0745$ while the heuristic fast
path carries detection, attributing the blindness to distribution transfer and
corroborating Sommer and Paxson~\cite{sommer2010}. Here we report a
\emph{second, independent} cause, obtained not by replay but by auditing the
compiled models and the feature-extraction code directly (DAY~255). It does not
replace the transfer explanation; it sits beneath it and is separately
reproducible.

\paragraph{Method.}
Every claim is measured against source and against the auto-generated forest
files. Two initial hypotheses were refuted by measurement before the finding was
reached, and are recorded for transparency: (i)~that the multi-flow aggregator
is never injected in production --- refuted, the initialization path is
unconditional and cannot fail on environment; (ii)~that
\texttt{source\_ip\_dispersion} is never computed --- refuted, it is computed
from the aggregator and the forest barely splits on it.

\paragraph{The sentinel and the tree.}
Features whose computation is unavailable return the shared constant
$\texttt{MISSING\_FEATURE\_SENTINEL} = -9999.0\text{f}$. Since every learned
threshold in the four forests lies in $[0,5.1]$, a served value of $-9999$
satisfies $\texttt{feature} \le \texttt{threshold}$ at every node that evaluates
it, forcing the left branch unconditionally. This is benign only if the feature
was also constant during training. It was not, and this is provable without the
training artifacts: \emph{a decision-tree split is selected only where a feature
has non-zero variance in the training set} (a constant column yields zero
information gain and is never chosen). Therefore the existence of even one split
node on a feature establishes that the feature varied at training time; served
as a constant at inference, it constitutes a train/serve distribution mismatch.
The forest was trained in a feature space the deployment does not reproduce.

\paragraph{The audit.}
For each head we verified (a)~that the deployed forest is the trained forest and
(b)~which feature index each split node decides on, cross-referenced against the
census of features returning the sentinel on every run. Index-to-name mappings
were verified against the auto-generated comments in each forest file and against
the detector's \texttt{to\_array()} serialization, not inferred. For the three
heads with a training-side counterpart (DDoS, Traffic, Internal), the deployed
forest is byte-identical to the trained one modulo a compiler-warning wrapper
(anonymous namespace), verified by \texttt{diff}: the served model \emph{is} the
validated model, so the defect lies entirely in the input vector.
Table~\ref{tab:deadfeatures} reports the result.

\begin{table}[h]
    \centering
    \footnotesize
    \caption{Split nodes deciding on absence-sentinel features, per head
        (DAY~255). ``Splits on sentinel'' counts decision nodes whose feature is
        served as $-9999.0\text{f}$ on every run; each routes left unconditionally.
        The Ransomware head is the extreme case: its highest-importance feature
        (\texttt{io\_intensity}, 24\%, recorded in the detector header) measures
        disk I/O, physically absent from network traffic.}
    \label{tab:deadfeatures}
    \begin{tabular}{lrrrl}
        \toprule
        \textbf{Head} & \textbf{Splits on sent.} & \textbf{Total splits} &
        \textbf{\%} & \textbf{Sentinel feature(s)} \\
        \midrule
        DDoS       & 16  & 256   & 6.25\% & geographical\_concentration \\
        Traffic    & 44  & 457   & 9.6\%  & tcp\_udp\_ratio, flow\_duration\_std, protocol\_variety \\
        Internal   & 21  & 420   & 5.0\%  & connection\_duration\_std \\
        Ransomware & 809 & 1{,}932 & 41.9\% & io\_intensity (24\% imp.), resource\_usage, file\_operations, process\_anomaly \\
        \midrule
        \textbf{Total} & \textbf{890} & \textbf{3{,}065} & \textbf{29.0\%} & --- \\
        \bottomrule
    \end{tabular}
\end{table}

\paragraph{The clean case: geography that is never observed.}
The DDoS head splits 16 times on \texttt{geographical\_concentration}. The
pipeline never captures or enriches with geolocation --- the feature was
deferred to a phase that was never implemented, and is served as the sentinel on
every run. By the argument above, the 16 splits prove the training data carried
real, separable geographic values. An \emph{optional} enrichment field, intended
to be resolved on demand for post-hoc context, was wired into a \emph{mandatory}
model input vector; the model has no notion of optionality and learned to depend
on a column the deployment leaves constant. This is a category error at the
train/serve boundary, not a weak feature.

\paragraph{The structural case: a host signal in a network sensor.}
The Ransomware head devotes 41.9\% of its splits to four features
(\texttt{io\_intensity}, \texttt{resource\_usage}, \texttt{file\_operations},
\texttt{process\_anomaly}) that measure disk, CPU, filesystem, and process
behavior. These are not unimplemented TODOs recoverable by more sniffer
engineering: they are host-domain signals that \emph{by construction do not
travel on the network}. The head's single most important feature (24\%) is among
them. No network-only model can serve them; the resolution is architectural, not
a bug fix (\S\ref{sec:eval:remediation}, \S\ref{sec:future:wazuh}).

\paragraph{Significance.}
The contribution is not the defect but the reproducible forensic record: a set of
classifiers with high in-sample metrics that, in deployment, decide nearly a third
of their splits on constants --- a failure mode provable from the compiled
artifacts alone, rarely audited because there is no incentive to, and, for
critical infrastructure, worse than no detector because it manufactures false
confidence. It is concurrent with, and independent of, the distribution-transfer
result of \S\ref{sec:eval:xgboost:ood} and \S\ref{sec:eval:bias}.


% ╔══════════════════════════════════════════════════════════════════════════╗
% ║ BLOQUE 3 — REMEDIATION PLAN (nueva subsección)                            ║
% ║ Insertar tras §sec:eval:deadfeatures                                       ║
% ╚══════════════════════════════════════════════════════════════════════════╝

\subsection{Remediation Plan: Per-Head Cost of Repair}
\label{sec:eval:remediation}

The audit partitions the four heads by a single physical criterion --- whether
the sentinel-served feature is computable from network traffic at all --- which
determines whether repair is a matter of implementation, retraining, or
architectural change. Table~\ref{tab:remediation} summarizes. The plan is
specified here; its measured outcomes are the subject of ongoing work and are not
claimed in this revision. Costs are conditional on a provenance measurement (the
training-time distribution of each column) that quantifies, but does not change,
the direction established above.

\begin{table}[h]
    \centering
    \footnotesize
    \caption{Per-head remediation. ``Computable in network?'' is the physical
    criterion that determines the repair class. The multi-flow aggregator on
    which Traffic and Internal depend is already live in production.}
    \label{tab:remediation}
    \begin{tabular}{lp{3.1cm}p{4.7cm}l}
        \toprule
        \textbf{Head} & \textbf{Computable in network?} & \textbf{Repair} & \textbf{Cost} \\
        \midrule
        Internal & Yes (aggregator, live) &
        Implement \texttt{connection\_duration\_std} from the aggregator; retrain if
        the served range diverges from training & Low \\
        \addlinespace
        Traffic & Partly (2/3 via aggregator; \texttt{tcp\_udp\_ratio} needs a
        protocol field) & Implement the three features; retrain as above & Low--Med \\
        \addlinespace
        DDoS & On demand, off the hot path (GeoIP latency) &
        Remove \texttt{geographical\_concentration} from the model input and retrain;
        keep it in the contract as an \emph{optional} enrichment field for post-hoc
        context (decouple model input from enrichment). Whether the repaired ML head
        outperforms the fast path's rate heuristics is measured, not assumed & Low \\
        \addlinespace
        Ransomware & \textbf{No} (host signals) &
        (A)~retrain network-only, losing the 24\% feature and 41.9\% of splits ---
        an honest but substantially weaker model; \textbf{(B)~source the four
        features from Wazuh host telemetry and fuse at the graph} --- the principled
        fix (\S\ref{sec:future:wazuh}) & (A)~Med / (B)~High \\
        \bottomrule
    \end{tabular}
\end{table}

\paragraph{Governing principle.}
The systemic fix is a single invariant on the model boundary: \emph{the model
input vector must contain only deterministically-present features; any feature
that can be absent is either excluded from the model or has its absence encoded
identically in training and serving}. The current design violates this by
training on real values and serving the sentinel. Three heads (Internal, Traffic,
DDoS) satisfy the invariant by implementing or removing the offending column and
retraining, and remain inline in the sub-microsecond detector. The Ransomware
head cannot satisfy it within the network sensor at all: it exits the detector to
the host domain (\S\ref{sec:future:wazuh}), where the response path becomes
asynchronous --- host evidence accumulates and, on high-confidence detection,
emits a block to the firewall \texttt{ipset} as a second evidence source
alongside the network detector. DDoS, by contrast, remains inline and real-time:
its signal is network-native and volumetric blocking cannot be deferred.


% ╔══════════════════════════════════════════════════════════════════════════╗
% ║ BLOQUE 4 — RANSOMWARE / HOST-DOMAIN                                        ║
% ║ 4A → añadir en §"Ransomware Detection: Scope and Validation"               ║
% ║ 4B → añadir en §future:wazuh (FEAT-INT-1)                                  ║
% ╚══════════════════════════════════════════════════════════════════════════╝

% ── [4A] tras "What is not validated." ──────────────────────────────────────
\paragraph{What the DAY~255 audit changes.}
The dead-feature audit (\S\ref{sec:eval:deadfeatures}) shows that the Ransomware
head's limitation is not a matter of training-data age but of \emph{sensor
domain}: four of its ten features --- including its highest-importance feature at
24\% --- measure disk I/O, CPU, filesystem, and process behavior, none of which
exist in a network capture. The behavioral indicators validated on Neris (SMB
burst, connection rate, DNS anomaly) are the network-observable shadow of
ransomware activity; the encryption behavior that \emph{defines} ransomware is a
host phenomenon. A network-only ransomware classifier therefore has a principled
ceiling, independent of retraining. The resolution is host-domain fusion
(\S\ref{sec:future:wazuh}).

% ── [4B] tras el párrafo existente de §future:wazuh ─────────────────────────
\paragraph{From enhancement to requirement (DAY~255).}
The audit in \S\ref{sec:eval:deadfeatures} sharpens the motivation for this
integration from correlation quality to necessity. The four features that cripple
the network-side Ransomware head --- disk I/O, resource usage, file operations,
process anomaly --- are precisely the signals a HIDS observes natively. Wazuh,
currently deployed to police the integrity of the pipeline's own components, is
the sensor whose domain contains exactly what the network sensor cannot produce.
The pipeline already carries the host ingestion path end-to-end (a Wazuh adapter
emitting a sealed \texttt{host\_domain\_v1} bronze record, a host bronze-to-gold
converter, and a loader into a dedicated host Kuzu graph), gated in the same
quality harness as the network path. What remains is not plumbing but two pieces:
a \emph{driver} that induces ransomware-characteristic host behavior on a
Wazuh-monitored, isolated VM so the agent emits the corresponding host telemetry,
and a host-domain model trained on that telemetry rather than on packet features.
The network and host graphs are then correlated at the
\texttt{community\_id}/\texttt{flow\_uid} substrate of \S\ref{sec:eval:parity}:
lateral-movement spread on the network and file encryption on the host become a
single composite signature that neither lens resolves alone. Because host
evidence is asynchronous, the host detector's response path is decoupled from the
inline pipeline: on high-confidence detection it emits a block to the firewall
\texttt{ipset} as a second evidence source, rather than through the
sub-microsecond scoring path. This re-derives, from a measured failure, the
danger-model stance --- detect damage in context rather than classify isolated
packets --- and is the primary architectural direction for the Ransomware head.

\paragraph{A note on training-data provenance for the host model.}
A natural question is whether the existing academic corpora (CTU-13, CIC-IDS) can
supply this host telemetry by routing them through Wazuh. They cannot: these are
network captures, and the host signals (disk, filesystem, process) were never
recorded in them --- the same category error the audit identifies, viewed from
the data side. Two viable sources exist instead: (i)~mapping an existing
\emph{host-behavior} ransomware dataset (system-call, FIM, or I/O traces) onto the
\texttt{host\_domain\_v1} contract; or (ii)~generating fresh host telemetry by
running the characteristic file-encryption I/O pattern --- via controlled benign
emulation, not live malware --- on the Wazuh-monitored VM already present in the
pipeline, and harvesting the agent output through the existing adapter. Either
path yields a host-domain corpus of the correct medium; neither is a retrofit of
a network dataset.